\documentclass{beamer}
\usepackage{graphicx} % Required for inserting images
\usepackage{german}
\usepackage{amsmath}
\usepackage{amssymb}
\usepackage{amsthm}
\usepackage{url}
\usepackage{xcolor} %vonmir

\usepackage{algorithm} %vonmir
\usepackage[noend]{algpseudocode} %vonmir; noend hilft das kein endif endfor am Ende dasteht


\newcommand{\Z}{\mathbb{Z}}% ganze Zahlen, fehlte in Vorlage
\newcommand{\N}{\mathbb{N}}% natürliche Zahlen, fehlte in Vorlage
\newcommand{\R}{\mathbb{R}}% reele Zahlen
\newcommand{\C}{\mathbb{C}}% komplexe Zahlen
\newcommand{\Q}{\mathbb{Q}}% rationale Zahlen
\newcommand{\beweis}[1]{\textbf{\structure{Beweis~#1.}}}


\newtheorem{thm}{Satz}
\newtheorem{cor}[thm]{Korollar}
\newtheorem{lem}[thm]{Lemma}
\newtheorem{prop}[thm]{Proposition}


\setbeamertemplate{theorems}[numbered] %von mir zum Nummerieren
%\setbeamertemplate{theorems}[normal font]

\AtBeginSection[]{
\begin{frame}{Überblick}
\tableofcontents[currentsection]
\end{frame}
}

\AtBeginSubsection[]{
\begin{frame}{Überblick}
\tableofcontents[currentsection, currentsubsection]
\end{frame}
} %von mir

\title{Proseminar: Primzahltest und Berechenbarkeit}
\author{Jana Chung}
\date{\today}

\begin{document}

% Title Slide
\begin{frame}
    \titlepage
\end{frame}

% First Content Slide
\begin{frame}{Inhaltsverzeichnis}
    \begin{enumerate}
        \item Was ist ein Algorithmus
        \item Berechnungsproblem
        \item Pseudocode
        \item Primzahltest
       \item Sieb des Eratosthenes
        \item Halteproblem
        
    \end{enumerate}
\end{frame}

\begin{frame}{1 Was ist ein Algorithmus}
    Ein Algorithmus ist eine endliche, wohldefinierte Folge von Anweisungen, die für
jede zulässige Eingabe nach endlich vielen Schritten terminiert und eine Ausgabe
liefert 
    \begin{itemize}
        \item  Welche Eingaben akzeptiert er?
        \item  Welche Rechenschritte werden in welcher Reihenfolge durchgeführt?
        \item  Wann stoppt der Algorithmus? Was wird dann ausgegeben?
    \end{itemize}

\end{frame}

\begin{frame}{2 Berechnungsprobleme}
    Durch eine Funktion $f\colon A \to B$ können wir ausdrücken dass ein Computerprogramm bei Eingabe eines Inputs \(a \in A\) das Ergebnis \(f(a) \in B\) ausgibt.

\begin{definition}   
Sei A eine nichtleere endliche Menge. Für $k \in \N \cup \{0\}$ bezeichnen wir mit $A^k$ die Menge der Funktionen $f: \{1,\dots,k\} \rightarrow A$.
Ein Element f $\in A^k$ heißt \textbf{Wort / Zeichenkette} der \textbf{Länge} k über dem \textbf{Alphabet} A. Wir schreiben es als Folge $f(1)... f(k)$.  
Das einzige Element $\emptyset$ von $A^0$ nennen wir das \textbf{leere Wort} (mit Länge 0). Wir setzen $A^\ast:= \bigcup_{k\in \N \cup \{0\}}^{} A^k$.
Eine \textbf{Sprache} über dem Alphabet A ist eine Teilmenge von $A^\ast$.
%Buch Def. 1.5.
\end{definition}

\end{frame}


\begin{frame}{2 Berechnungsprobleme}

\begin{definition}
     Ein \textbf{Berechnungsproblem} ist eine Relation $P \subseteq D \times E$, wobei es zu jedem d $\in$ D mindestens ein e $\in$ E gibt mit $(d,e) \in $ P. Wenn $(d,e) \in$ P, dann ist e eine \textbf{korrekte Ausgabe} für das Problem P mit Eingabe d. Die Elemente von D heißen \textbf{Instanzen} des Problems. P heißt \textbf{eindeutig}, wenn P eine Funktion ist. Sind D, E Sprachen über einem endlichen Alphabet A, so sprechen wir von einem \textbf{diskreten Berechnungsproblem}. Sind D und E Teilmengen von  $\R^m$ bzw. $\R^n$ für $m, n \in \N $, so sprechen wir von einem \textbf{numerischen Berechnungsproblem}. Ein eindeutiges Berechnungsproblem $P \colon D \to E$
     mit $\left| E \right| = 2$ heißt \textbf{Entscheidungsproblem}. 
     % Buch Def. 1.6
\end{definition}

\end{frame}

\begin{frame}{3 Pseudocode}
    Algorithmen sollen Berechnungsprobleme lösen, heutzutage erfolgt die Ausführung mit Computern.
    
    \begin{itemize}
        \item Algorithmus entwerfen
        \item wird in Computerprogramm (z.B. in C++) implementiert
        \item wird zu Maschinencode kompiliert
        \item Ausführung
    \end{itemize}


 \textbf{Pseudocode}: macht einen Algorithmus konkret in einer noch für Menschen verständlichen Sprache, und gleichzeitig leicht in eine Programmiersprache umsetzbar. Gut um Idee des Algorithmus darzustellen   
\end{frame}

\begin{frame}{4 Primzahltest - Was heißt prim?}

 Der Algorithmus soll herausfinden, ob eine gegebene natürliche Zahl \textbf{prim ist oder nicht}. 

 \vspace{0.5em}

   Entscheidungsproblem: \(
   \{(n,e) \in \N \times \{ 0, 1\} | \, e= 1 \iff
   \) n prim $\}$.
   
   \begin{definition}
       Eine natürliche Zahl $n$ heißt prim, wenn \(
       n \geq 2
       \) ist und es keine natürliche Zahl $a$ und $b$ gibt mit \( a > 1, b > 1\) und $n= a \cdot b$
   \end{definition}

  
   
    %\begin{thm}
     %   First theorem
    %\end{thm}
\end{frame}

\begin{frame}{4 Primzahltest - Code und Definition}
    
    \begin{itemize}
        \item \textbf{If-Anweisung:} Führt Anweisungen abhängig vom Wahrheitswert einer Bedingung aus. Ist die Bedingung wahr, wird der \texttt{then}-Teil ausgeführt, andernfalls der \texttt{else}-Teil (optional)
        \item \textbf{For-Anweisung:} Eine Variable durchläuft sukzessive Werte, und für jeden Wert wird ein Anweisungsblock ausgeführt.
    \end{itemize}

\begin{definition} [Gaußklammern]
Für \(
x \in \R
\) definieren wir:
\(
\lfloor x \rfloor := \max\{k \in \Z | k \leq x \}, \quad
\lceil x \rceil := \min \{k \in \Z | k \geq x \}.
\)
\end{definition}

\begin{definition}[Modulo-Operation]
Für $a,b \in \N$ definieren wir:
\(
a \bmod b := a - b \cdot \lfloor \frac{a}{b} \rfloor .
\)
Es gilt
\(
b | a \iff a \bmod b = 0.
\)
    
\end{definition}
\end{frame}

\begin{frame}{4 Primzahltest - Algorithmus}
   
   Nun erfolgt der Algorithmus, ob eine natürliche Zahl prim ist oder nicht in Pseudocode
   
\begin{algorithm}[H]
\begin{algorithmic}[1]
\Require $n \in \mathbb{N}$
\Ensure Antwort auf die Frage, ob $n$ prim ist
\If{$n = 1$}
    \State $\text{antwort} \gets \text{``nein''}$
\Else
    \State $\text{antwort} \gets \text{``ja''}$
    \For{$i \gets 2$ \textbf{to} $\lfloor \sqrt{n} \rfloor$}
        \If{$i \mid n$}
            \State $\text{antwort} \gets \text{``nein''}$
        \EndIf
    \EndFor
\EndIf
\State \Return $\text{antwort}$
\end{algorithmic}
\end{algorithm}

\end{frame}


\section{Eratosthenes}
\begin{frame}{5 Sieb des Eratosthenes – Überblick}
\begin{itemize}
   \item Algorithmus
    \item Beweis
\end{itemize}
\end{frame}

\subsection{Algorithmus}
\begin{frame}{5 Sieb des Eratosthenes - Algorithmus}

Nach Eingabe einer natürlichen Zahl n besteht die Aufgabe darin, alle Primzahlen p mit p $\le$ n zu berechnen. Der Algorithmus lautet wie folgt:

\begin{algorithm}[H]
\begin{algorithmic}[1]
\Require $n \in \mathbb{N}$
\Ensure alle Primzahlen $\leq n$

\For{$i \gets 2$ \textbf{to} $n$}
    \State $p[i] \gets \text{``ja''}$
\EndFor

\For{$i \gets 2$ \textbf{to} $n$}
    \If{$p[i] = \text{``ja''}$}
        \For{$j \gets i$ \textbf{to} $\lfloor \frac{n}{i} \rfloor$}
            \State $p[i \cdot j] \gets \text{``nein''}$
        \EndFor
    \EndIf
\EndFor

\State \Return alle $i$ mit $p[i] = \text{``ja''}$
\end{algorithmic}
\end{algorithm}
    
\end{frame}

\begin{frame}[allowframebreaks]{5 Sieb des Eratosthenes}
\begin{thm}
   Obiger Algorithmus ist korrekt.
\end{thm}

Wir zeigen die Korrektheit des Algorithmus. Dazu beweisen wir für alle 
\( k \in \{2, \dots, n\} \):

\textbf{\[
k \text{ wird genau dann ausgegeben, wenn } k \text{ prim ist.}
\]}
Dies beweisen wir indem wir folgende 2 Aussagen beweisen

\begin{itemize}
    \item Ist k prim, so wird k ausgegeben (1)
    \item Ist k nicht prim, so wird k nicht ausgegeben (2)
\end{itemize}
\medskip
\end{frame}

\subsection{Beweis}
\begin{frame}{5 Sieb des Eratosthenes}

\textcolor{structure}{\textbf{Beweis 1.}}

    \textbf{(1) Ist \(k\) prim, so wird \(k\) ausgegeben.}

Sei \(k \in \{2, \dots, n\}\) eine Primzahl. Im Algorithmus wird eine Zahl \(k\) genau dann nicht ausgegeben, wenn irgendwann 
\(p[k] \leftarrow \text{``nein''}\) gesetzt wird. Wir zeigen, dass dies für Primzahlen nicht passiert.

Die Zuweisung \(p[k] \leftarrow \text{``nein''}\) erfolgt nur dann, wenn es Zahlen 
\(i, j \in \mathbb{N}\) mit \(i \ge 2\), \(j \ge i\) gibt, sodass \[k = i \cdot j.\]
Das bedeutet aber, dass \(k\) als Produkt zweier Zahlen größer oder gleich 2 dargestellt werden kann. Dies widerspricht der Annahme, dass \(k\) prim ist.
Also wird für eine Primzahl \(k\) niemals \(p[k] = \text{``nein''}\) gesetzt. Somit gilt während des gesamten Algorithmus: \[p[k] = \text{``ja''}.\] Folglich wird \(k\) ausgegeben.
\end{frame}


\begin{frame}{5 Sieb des Eratosthenes}

 \textbf{(2) Ist \(k\) nicht prim, so wird \(k\) nicht ausgegeben.}

Sei nun \(k \in \{2, \dots, n\}\) keine Primzahl. Dann ist \(k\) zusammengesetzt, 
also existieren Zahlen \(i, j \in \mathbb{N}\) mit
\[
k = i \cdot j \quad \text{und} \quad 2 \le i \le j.
\]

Wir wählen \(i\) als den kleinsten Teiler von \(k\), der mindestens 2 ist. 
Dann gilt insbesondere:
\[
2 \le i \le j = \frac{k}{i}.
\]

Da \(i\) der kleinste Teiler von \(k\) ist, kann \(i\) selbst keine nichttrivialen 
Teiler besitzen. Also ist \(i\) eine Primzahl.

Nach Teil (1) wissen wir, dass für Primzahlen stets \(p[i] = \text{``ja''}\) gilt, 
wenn der Algorithmus die Zahl \(i\) betrachtet.
    
\end{frame}


\begin{frame}{5 Sieb des Eratosthenes}
In diesem Schritt führt der Algorithmus die innere Schleife aus und setzt für alle 
geeigneten \(j\):
\[
p[i \cdot j] \leftarrow \text{``nein''}.
\]
Da \(k = i \cdot j\), wird insbesondere auch
\[
p[k] \leftarrow \text{``nein''}
\]
gesetzt.
Ab diesem Zeitpunkt bleibt \(p[k] = \text{``nein''}\), da dieser Wert im Algorithmus 
nicht mehr geändert wird.

Wenn später im äußeren Durchlauf der Wert \(i = k\) erreicht wird, gilt daher:
\[
p[k] = \text{``nein''},
\]
und somit wird \(k\) nicht ausgegeben.

\medskip

Damit ist die Korrektheit des Algorithmus gezeigt.
\hfill$\square$

\end{frame}

\section{Halteproblem}
\begin{frame}{6 Halteproblem – Überblick}
\begin{itemize}
    \item Kodierung von Wörtern (Lemma)
    \item Abzählbarkeit von Programmen
    \item Existenz nicht berechenbarer Funktionen
    \item Definition der Haltefunktion
    \item Widerspruchsbeweis
\end{itemize}
\end{frame}

   \subsection{Kodierung von Wörtern}
\begin{frame}{6 Halteproblem}

    Zunächst benötigen wir einige Vorarbeit.

    \begin{lem}
    Seien $k,l \in \mathbb{N}$. Definiere die Abbildung
    \[ 
    f^l: \{0,\dots,k-1\}^l \to \{0,\dots,k^l-1\}, \quad
    f^l(w) := \sum_{i=1}^{l} a_i k^{\,l-i}.
    \]
   für alle \( w = a_1,\dots,a_l \in \{0,...,k-1\}^l \).
    Dann ist die Abbildung $f^l$ wohldefiniert und bijektiv.
\end{lem}
\end{frame}

\begin{frame}{6  Halteproblem}

%\begin{block}{Beweis X}
\beweis{2}

    \textbf{1. Wohldefiniertheit:}

Sei $(a_1,\dots,a_l) \in \{0,\dots,k-1\}^l$. Dann gilt für alle $i$:
\[
0 \le a_i \le k-1.
\]
Damit folgt
\[
0 \le a_i k^{l-i} \le (k-1)k^{l-i}.
\]
Durch Aufsummieren erhalten wir
\[
0 \le f_l(a_1,\dots,a_l) = \sum_{i=1}^l a_i k^{l-i}
\le \sum_{i=1}^l (k-1)k^{l-i}.
\]
Die rechte Summe ist eine geometrische Reihe:

\end{frame}

\begin{frame}{6 Halteproblem}
Wir formen die rechte Summe um:
\[
\sum_{i=1}^l (k-1)k^{l-i}
= (k-1)\sum_{i=1}^l{k^{l-i}}
= (k-1)\sum_{i=0}^{l-1}{k^{l-i-1}} \]
\[
= (k-1)\sum_{j=0}^{l-1} k^j
= (k-1)\cdot \frac{k^l - 1}{k-1}
= k^l - 1.
\]
    Also gilt insgesamt
\[
0 \le f_l(a_1,\dots,a_l) \le k^l - 1,
\]
und damit ist $f_l$ wohldefiniert.
\end{frame}

\begin{frame}{6 Halteproblem}
    \textbf{2. Injektivität:}

Seien $(a_1,\dots,a_l),(b_1,\dots,b_l) \in \{0,\dots,k-1\}^l$ mit
\[
(a_1,\dots,a_l) \ne (b_1,\dots,b_l)
\] 
beliebig aber fest gewählt.
Dann gibt es einen kleinsten Index $j \in \{1,\dots,l\}$ mit
\[
a_j \ne b_j.
\]
Ohne Einschränkung sei $a_j > b_j$.

\[
f^l(a_1,\dots,a_l) - f^l(b_1,\dots,b_l)
= \sum_{i=1}^l (a_i - b_i)k^{l-i}.
\]
Für $i<j$ gilt $a_i = b_i$, also verschwinden diese Terme:
\[
= (a_j - b_j)k^{l-j} + \sum_{i=j+1}^l (a_i - b_i)k^{l-i}.
\]
\end{frame}

\begin{frame}{6 Halteproblem}
    Wir betrachten die Differenz:


Nun schätzen wir die Summe für $i>j$ nach unten ab. Da $a_i,b_i \in \{0,\dots,k-1\}$ gilt
\[
a_i - b_i \ge -(k-1).
\]
Also
\[
\sum_{i=j+1}^l (a_i - b_i)k^{l-i}
\ge -\sum_{i=j+1}^l (k-1)k^{l-i}.
\]
Wie oben ergibt die geometrische Summe:
\[
\sum_{i=j+1}^l (k-1)k^{l-i}
= k^{l-j} - 1.
\]
Damit folgt
\[
\sum_{i=j+1}^l (a_i - b_i)k^{l-i}
\ge -(k^{l-j} - 1).
\]

\end{frame}

\begin{frame}{6 Halteproblem}
    
Zusammen ergibt sich:
\[
f^l(a_1,\dots,a_l) - f^l(b_1,\dots,b_l)
\ge (a_j - b_j)k^{l-j} - (k^{l-j} - 1).
\]
Da $a_j > b_j$, gilt $a_j - b_j \ge 1$, also:
\[
\ge k^{l-j} - (k^{l-j} - 1) = 1.
\]
Insbesondere ist
\[
f^l(a_1,\dots,a_l) > f^l(b_1,\dots,b_l).
\]
Also ist $f^l$ injektiv.

\end{frame}

\begin{frame}{6 Halteproblem}

\textbf{3. Surjektivität:}

Die Mengen haben dieselbe Mächtigkeit:
\[
|\{0,\dots,k-1\}^l| = k^l = |\{0,\dots,k^l-1\}|.
\]

Die linke Menge besteht aus allen Tupeln der Länge l. Für jede Position gibt es k verschiedene Möglichkeiten. Die Anzahl ergibt sich also als Produkt $k\cdot k\cdot k \dots k = k^l$. Die rechte Menge besteht einfach aus allen Zahlen von 0 bis $k^l-1$, weist also auch die Anzahl $k^l$ auf (am Ende zieht man zwar 1 ab, fängt aber mit 0 an).

\end{frame}

\begin{frame}{6 Halteproblem}

Da $f^l$ eine injektive Abbildung zwischen zwei endlichen Mengen gleicher Größe ist, ist sie automatisch auch surjektiv.

\vspace{0.5em}

$f^l$ ist somit bijektiv. %$\qed$
\hfill$\square$
    
\end{frame}

\subsection{Abzählbarkeit}
\begin{frame}{6 Halteproblem}

    \begin{thm}
    Sei A eine nichtleere Menge. Dann ist die Menge $A^*$ abzählbar.
\end{thm}
\vspace{0.5em}
Der Beweis wird in Folge nur skizziert:
\end{frame}

\begin{frame}{6 Halteproblem}

    

\beweis{3}
Die Idee des Beweises besteht darin, die Wörter aus $A^*$ systematisch zu nummerieren.

Zunächst werden alle Wörter nach ihrer Länge geordnet. Für jedes $l \in \N$ gibt es genau $|A|^l$ Wörter der Länge $l$. Da $A$ endlich ist, ist somit auch jede Menge $A^l$ endlich.

Anschließend werden die Wörter fester Länge $l$ eindeutig als Zahlen im $|A|$-adischen Zahlensystem kodiert. Dazu wird jedem Symbol aus $A$ eine Ziffer aus $\{0,\dots,|A|-1\}$ zugeordnet.

Durch diese Kodierung erhält man eine injektive Abbildung von $A^*$ nach $\N$, da Wörter unterschiedlicher Länge in verschiedene Bereiche fallen und die Darstellung innerhalb fester Länge eindeutig ist.

Damit ist $A^*$ abzählbar.
\hfill$\square$
\end{frame}

\begin{frame}{6 Halteproblem}

\begin{cor}
\label{abzählbar}
    Die Menge aller C++-Programme ist abzählbar.
\end{cor}

\beweis{4}
    C++-Programme sind Wörter über einem endlichen Alphabet. Die zu zeigende Aussage folgt also aus Satz \ref{abzählbar}.  
    \hfill$\square$
\end{frame}

\begin{frame}{6 Halteproblem}
\begin{thm}
    Es gibt Funktionen \(f: \N\to \{0,1\} \), die von keinem C++-Programm berechnet werden.

\end{thm}

\beweis{5}
    Sei $\mathcal{P}$ die nach Korrolar \ref{abzählbar} abzählbare Menge der C++-Programme, die eine Funktion \(f: \N\to \{0,1\} \) berechnen.

    Sei \(g: \mathcal{P}\to \{0,1\} \) eine injektive Funktion. 
    Für \(P \in \mathcal{P}\) sei $f^P$ die von P berechnete Funktion.

    Betrachte eine Funktion \(f: \N\to \{0,1\} \) mit \( f(g(P)):=1-f^P(g(P))\) für alle \(P \in \mathcal{P}\).
    Eine solche existiert, da g injektiv ist. Offenbar ist \(
   f \neq f^P 
    \)
    für alle 
    \(
    P \in \mathcal{P}
    \), und somit wird f von keinem C++ Programm berechnet.
    \hfill$\square$
\end{frame}

\subsection{Halteproblem}
\begin{frame}{6 Halteproblem}
Unter den nicht berechenbaren Funktionen geht es nun um das berühmteste Beispiel, dem \textbf{Halteproblem}: Es gibt kein Programm, dass für alle Programme und Eingaben entscheiden kann, ob das Programm terminiert.

\medskip

Wir betrachten die natürlichen Zahlen $\N$ und die Menge aller Programme $Q$. Da Programme endliche Texte darstellen, können wir alle Programme durchnummerieren. Es gibt also eine surjektive Funktion \(g: \N \to Q \) - jede Zahl kodiert ein Programm.
Genauso geht die Kodierung auch in die andere Richtung, jedes Programm kodiert eine Zahl: \( f: Q \to \N
\). Jedem Programm wird eine Zahl zugeordnet.
\end{frame}

\begin{frame}{6 Halteproblem}
Wir definieren die Funktion wie folgt:
\begin{align}
    h: \N \times \N \to \{1,0\}
\end{align}

\begin{align}
    h(x,y) = \begin{cases}
        1,& \text{wenn Programm g(x) bei Eingabe y terminiert}\\
        0,& \text{sonst}
    \end{cases}
\end{align}

x: Nummer eines Programms; y: Eingabe; h(x,y) sagt ob Programm terminiert oder nicht.

\medskip

Diese Funktion heißt \textbf{Haltefunktion} und ihr zugehöriges Entscheidungsproblem heißt \textbf{Halteproblem}.
\end{frame}

\subsection{Widerspruchsbeweis}
\begin{frame}{6 Halteproblem}
\begin{thm}
    Die Haltefunktion wird von keinem C++ Programm berechnet. Das heißt es existiert kein Programm das $h$ berechnet.
\end{thm}

\beweis{6}
    Die Beweisidee ist ein Widerspruchsbeweis.
    
    Angenommen, es existiert ein Programm $P$, das $h(x,y)$ korrekt berechnet. Wir konstruieren ein neues Programm $Q$, das wie folgt funktioniert:

    Bei Eingabe $x \in \N$:
    \begin{enumerate}
        \item Berechne h(x,x)
        \item Falls \(h(x,x) = 0 \rightarrow \) terminiert
        \item falls \(h(x,x) = 1\rightarrow \) terminiert nicht 
    \end{enumerate}
    
   
    \begin{block}{Schlüsselidee}
Es passiert genau das \alert{Gegenteil} von dem, was \( h \) vorhersagt.
\end{block}

\end{frame}

\begin{frame}{6 Halteproblem}


   Da $Q$ ein Programm ist und jedes Programm eine Nummer hat, gibt es ein $q$ mit \( g(q) = Q\).

    Nun betrachten wir $h(q,q)$.
    \begin{itemize}
        \item \textbf{Fall 1:} Betrachte \(h(q,q) = 1\). 
    $Q(q)$ terminiert, aber nach Definition von $Q$ terminiert es nicht.
         \item \textbf{Fall 2:} Betrachte \(h(q,q) = 0\).
    $Q(q)$ terminiert es nicht, aber nach Definition von $Q$ müsste es terminieren.

    \end{itemize}
    
\medskip

    In beiden Fällen entsteht also ein Widerspruch. Also war die Annahme falsch. Die Haltefunktion ist somit nicht berechenbar.
    \vspace{0.5em}

Man sagt dazu: das Halteproblem ist \textbf{nicht entscheidbar}. \hfill$\square$
\end{frame}

\begin{frame}{6 Halteproblem}
Ein gutes Beispiel für ein Programm, bei dem nicht leicht zu sehen ist, ob es hält, ist die \textbf{Collatz-Folge}. Der Algorithmus funktioniert wie folgt:
Nach Eingabe von \(n \in \N \) wird berechnet, ob die Collatz-Folge von n den Wert 1 erreicht. Dabei geht das Programm wie folgt vor. Wenn n eine gerade Zahl ist, so wird n halbiert. Ansonsten wird es verdreifacht und um 1 erhöht. Dies wiederholt sich bis die Collatz-Folge den Wert 1 erreicht. Bis heute konnte nicht erfasst werden, ob dieser Algorithmus immer, also für jedes n, hält.

\end{frame}

\begin{frame}{Literatur}

\textbf{Hauptquelle:}

\begin{itemize}
    \item B. H. Hougardy: \\
    \textit{Algorithmische Mathematik} \\
    Springer, 2010
\end{itemize}

\vspace{1em}

\textbf{Bemerkung:}
\begin{itemize}
    \item Grundlage für alle Definitionen und Beweise
    \item Einige Beweisschritte wurden im Vortrag ausführlicher dargestellt
\end{itemize}

\end{frame}




\end{document}


%Useful: \hfill$\square$ für qed kästchen